\documentclass{report}
\usepackage{graphicx}
\usepackage{amssymb}
\usepackage{amsmath}
\usepackage{amsthm}

\theoremstyle{plain}
\newtheorem{theorem}{Theorem}[chapter] % reset theorem numbering for each chapter

\theoremstyle{definition}
\newtheorem{definition}[theorem]{Definition} % definition numbers are dependent on theorem numbers
\newtheorem{example}[theorem]{Example} % same for example numbers
\newtheorem{remark}[theorem]{Remark} % same for remarks

\begin{document}

\title{Notes on Analysis\\
	\Large Compiled from lectures by Prof. Vittal Rao\\
	\large IISc, Bangalore, IN}
\author{Soumyajit De}

\maketitle
%\tableofcontents
%\newpage
\chapter{Lectures 1}
\section{Functions}
\begin{definition}[Function]
\end{definition}
\subsection{Terminology}
\subsection{Function Extention}
\subsection{Inverse Image of a Function}
\subsection{Inverse Function}

\chapter{Lectures 17-19}
\section{Metric Space}
For the set $\mathbb{R}$, the operations (binary functions) `$+$' and `$\cdot$' provide the algebraic structure of a field, represented as the tuple $(\mathbb{R}, +, \cdot)$. However, to introduce a \emph{topological} structure for $\mathbb{R}$ - which is essential in order to define limit and convergence - we need to take into account of a different binary function, called the \emph{metric}. The set, along with the function, defines a metric space.
\begin{definition}[Metric Space]
A metric space for a set $S$ is defined as the tuple $(S, d)$ where $d:S\times S\mapsto[0, \infty)$ where the following conditions hold $\forall x, y, z\in S$:
\begin{itemize}
\item[(a)] $d(x,y)=d(y,x)$ [Symmetry]
\item[(b)] $d(x,y)=0\Leftrightarrow x=y$ [Uniqueness]
\item[(c)] $d(x,y)\le d(x,z)+d(y,z)$ [Triangular Inequality]
\end{itemize}
\end{definition}
\begin{example}
\noindent $(\mathbb{R},d)$, where $d(x,y)=|x-y|, \forall x,y\in\mathbb{R}$, where $|\cdot|:\mathbb{R}\mapsto[0,\infty)$ is the absolute-value function.
\end{example}

\section{Convergence}
\noindent With the notion of a metric, we can talk about \emph{how far} the elements are from one another. This notion of a \emph{distance} is essential to define the properties like convergence and limit. Note: a sequence $\{x_n\}$ ($\forall x_n\in S$) is countably infinite by definition, \emph{i.e.}, $n\in\mathbb{N}$.
\begin{definition}[Convergence and Limit]
The sequence $\{x_n\}$ is said to be converging towards $x_0\in S$ iff $\forall\epsilon>0$, $\exists N_\epsilon\in\mathbb{N}$, such that $\forall n >N_\epsilon$, $d(x_n, x_0) < \epsilon$. We denote this as $\{x_n\}\rightarrow x_0$ and we call $x_0$ the limit of the sequence $\{x_n\}$.
\end{definition}
\noindent This definition of ordinary convergence requires that the limit point $x_0$ is defined and is has to be present in the set $S$. There is another notion of convergence, called the Cauchy-convergence, which removes this requirement.
\begin{definition}[Cauchy Convergence]
The sequence $\{x_n\}$ is called Cauchy-convergent iff $\forall\epsilon>0$, $\exists N_\epsilon\in\mathbb{N}$, such that $\forall n,m >N_\epsilon$, $d(x_n, x_m) < \epsilon$.
\end{definition}
\begin{remark}
The beauty of convergence is that it preserves the algebraic operations defined on the set. More precisely, let $\{x_n\}\rightarrow x_0$ and $\{y_n\}\rightarrow y_0$ are two sequences defined in $\mathbb{R}$. Then the following are true:
\begin{itemize}
\item[(a)] $\{x_n+y_n\}\rightarrow x_0+y_0$
\item[(b)] $\{x_n\cdot y_n\}\rightarrow x_0\cdot y_0$
\end{itemize}
\end{remark}
\section{Open and Closed sets}
\noindent With the definition and properties of metric space, we can now define open and closed set. We begin by defining an $\epsilon$ \emph{open-ball}, $B_\epsilon(x)$ for a metric space $(S, d)$.
\subsection{Terminology}
\begin{definition}[Open-Ball]
$B_\epsilon(x):=\{y : d(x,y)<\epsilon\}$ for some $\epsilon>0$, $x,y\in S$.
\end{definition}
\begin{definition}[Internal Point]
$x\in S$ is called an internal point of $S$ iff $\exists \epsilon>0$, such that $B_\epsilon(x)\subseteq S$.
\end{definition}
\begin{definition}[Internal of a Set]
The set of all internal points of a given set $S$ is called the internal of a set, i.e. $S^0:=\{x : x\in S, \exists \epsilon>0, B_\epsilon(x)\subseteq S\}$.
\end{definition}
\noindent In general, it is easy to see that $S^0\subseteq S$. Using the definition of \emph{internal}, we can now define open-sets as follows.
\begin{definition}[Open Set]
A set $S$ is called open set iff $S=S^0$.
\end{definition}
\begin{definition}[Closed Set]
A set $S$ is called a closed set iff $S^c$ is open.
\end{definition}
\noindent Interesting note here: open and closed sets do not form a dichotomy. There are sets that are both open and closed (for example, the $\phi$ and $\mathbb{R}$ sets are both open and closed) and there are sets that are neither.
\subsection{Properties of Open Set}
\noindent Now we discuss the consequences of applying set operations on open sets. We can arrive at analogous (dual) properties for closed sets as well.
\begin{theorem}
Open sets are closed under union. Let $\{A_j\}_{j\in\mathcal{J}}$ be a collection of open sets (for any index set $\mathcal{J}$), Then $\bigcup_{j\in\mathcal{J}} A_j$ is also open.
\end{theorem}
\begin{proof}
Let $A_i$ and $A_j$ are two open sets. Now, $\forall z\in A_i\cup A_j\implies z\in A_i\vee z\in A_j$. Let's consider both these possibilities one by one. $z\in A_i\implies \exists \epsilon_i>0, B_{\epsilon_i}(z)\subseteq A_i\subseteq A_i\cup A_j$. Similarly, $z\in A_j\implies \exists \epsilon_j>0, B_{\epsilon_j}(z)\subseteq A_j\subseteq A_i\cup A_j$. So, in both these cases, $\exists\epsilon>0$ such that $\forall z\in A_i\cup A_j, B_\epsilon(z)\subseteq A_i\cup A_j$.

We can extend the same argument and conclude that union of any collection, finite, countable or uncountable, of open sets is an open set.
\end{proof}
\begin{theorem}
Open sets are closed under finite intersection. Let $\{A_i\}_{i=1}^N$ be a finite collection of open sets. Then the set $\bigcap_{i=1}^N A_i$ is also open.
\end{theorem}
\begin{proof}
Let $A_i$ and $A_j$ are two open sets. By definition, this means that $\exists \epsilon_i>0$ such that $\forall x\in A_i$, $B_{\epsilon_i}(x)\subseteq A_i$ and $\exists \epsilon_j>0$ such that $\forall y\in A_j$, $B_{\epsilon_j}(y)\subseteq A_j$. Now, $\forall z\in A_i\cap A_j\implies z\in A_i\wedge z\in A_j$. Let $\epsilon=\min(\epsilon_i, \epsilon_j)$. Since both $\epsilon_i>0$ and $\epsilon_j>0$, we know that $\min(\epsilon_i, \epsilon_j)>0$. Also, since $\epsilon<\epsilon_i$, we have $B_\epsilon(z)\subseteq B_{\epsilon_i}(z)\subseteq A_i$. Similarly, we have $B_\epsilon(z)\subseteq B_{\epsilon_j}(z)\subseteq A_j$. Combining these two, we have $\forall z\in A_i\cap A_j$, $\exists\epsilon>0$, such that $B_\epsilon(z)\subseteq A_i\wedge B_\epsilon(z)\subseteq A_j\implies B_\epsilon(z)\subseteq A_i\cap A_j$. This means, that $A_i\cap A_j$ is an open set.

We can extend this argument for any finite collection of sets, since for any finite collection of $\{\epsilon_1, \epsilon_2, \cdots \epsilon_n\}$, $\min_{i=1}^n \epsilon_i > 0$. However, this is not true for countably infinite sequences in general. For example, if we have a sequence $\{\epsilon_n\}$ as $\{1,\frac{1}{2},\frac{1}{4},\cdots\}$, then $\min_{n\in\mathbb{N}} x_n = 0$. But we need an $\epsilon>0$ to form an open set. Therefore, only in the finite case, the intersection is guaranteed to be an open set.
\end{proof}
\begin{example} The open interval, defined as $(a, b):=\{x : a < x < b\}$ is the simplest possible open set defined on $\mathbb{R}$. Similarly, the closed interval, defined as $[a, b]:=\{x : a \leq x \leq b\}$ is the simplest possible closed set in $\mathbb{R}$.
\end{example}
\subsection{Simple Construction of Sets}
Let's take the collection of all \emph{open balls} defined on a metric space $S$, called $\mathcal{O}$. We can take elements from this collection and start applying (countable or uncountable) union and finite intersection. The result sets are all open sets. Any possible open set in the metric space can be built in this fashion. The collection of all \emph{open sets} is usually called $\mathcal{G}$, and collection of all \emph{closed sets} is called $\mathcal{F}$.

In fact, it is even possible to construct $\mathcal{G}$ from a countable \emph{disjoint} set of open-balls from $\mathcal{O}$ and applying union. This is known as simple construction of sets (i.e. sets formed by any countable operations on other sets). This result is known as Lindeloff's theorem.

Interestingly, if we take a countable collection from $\mathcal{G}$ and apply \emph{intersection} instead, then we may or may not end up with an open set (since open sets are only guaranteed to be closed under \emph{finite} intersection). These new sets are called $\mathcal{G}_\delta$ sets, where the $\delta$ corresponds to countable intersections. Similarly, a countable union from $\mathcal{F}$ forms $\mathcal{F}_\sigma$ sets, $\sigma$ corresponding to the operations of performing countable unions.

We can now take these $\mathcal{G}_\delta$ and $\mathcal{F}_\sigma$ sets and apply further $\sigma$ and $\delta$ operations to construct $\mathcal{G}_{\delta\sigma}$, $\mathcal{F}_{\sigma\delta}$, $\mathcal{G}_{\delta\sigma\delta}$, etc.
\section{Measure Theory}
Let's say, we have a \emph{non-empty} set, $\Omega$. Clearly, the number of subsets are $2^{|\Omega|}$. We are interested a non-empty collection of subsets, $B$, which is \emph{self-contained}:
\begin{itemize}
\item[(1)] $x\in B\implies x^c\in B$ [closed under complement]
\item[(2)] $\{x_n\}\in B\implies \bigcup_{n\in\mathbb{N}} x_n \in B$ [closed under countable union]
\end{itemize}
\noindent We call this collection $B$ a $\sigma$-algebra of $S$.
\begin{remark}
$\sigma$-algebras are closed under countable intersection.
\end{remark}
\begin{proof}
\begin{align}
\{x_n\}\in B
&\implies \{x_n^c\}\in B &&\text{(from 1)} \\
&\implies \cup_{n\in\mathbb{N}} x_n^c \in B &&\text{(from 2)}\\
&\implies (\cup_{n\in\mathbb{N}} x_n^c )^c\in B &&\text{(from 1)} \\
&\implies \cap_{n\in\mathbb{N}} x_n \in B &&\text{(from De-Morgan's law)}&&\qedhere
\end{align}
\end{proof}
\begin{theorem}
Any $\sigma$-algebra for a set, $\Omega$, contains at least $\phi$ and $\Omega$.
\end{theorem}
\begin{proof}
Let $B$ be a $\sigma$-algebra of $\Omega$. Since $B$ is non-empty, we have at least an element $x\in B$.
\begin{align}
x\in B
&\implies x^c\in B \\
&\implies x\cap x^c=\phi\in B \\
&\implies x\cup x^c=\Omega\in B &&\qedhere
\end{align}
\end{proof}
\noindent So, for any non-empty set $\Omega$, there is a trivial $\sigma$-algebra, which contains just $\phi$ and $\Omega$. The collection of all subsets, $2^\Omega$, is also a $\sigma$-algebra, since it contains all possible subsets of $\Omega$, along with their complements and countable unions, including $\phi$ and $\Omega$.
\begin{theorem}
Intersection of $\sigma$-algebras for a set is a $\sigma$-algebra.
\end{theorem}
\begin{proof}
Let $B_i$ and $B_j$ be two $\sigma$-algebras of $\Omega$ and $B_i\cap B_j=B$. Since $\phi$ and $\Omega$ are present in both $B_i$ and $B_j$ (from theorem 3), they must be present in $B$. Therefore, $B$ is non-empty. Since it is non-empty, $\exists x\in B$. This implies that $x\in B_i\wedge x\in B_j$, which means, $x^c\in B_i\wedge x^c\in B_j$. Therefore, $x^c\in B_i\cap B_j=B$. Similarly, $\{x_n\}\in B\implies \{x_n\}\in B_i\wedge\{x_n\}\in B_j$. This concludes that $\cup_{n\in\mathbb{N}} x_n\in B_i\wedge \cup_{n\in\mathbb{N}} x_n\in B_j\implies \cup_{n\in\mathbb{N}} x_n\in B_i\cap B_j=B$. Therefore, $B$ is a $\sigma$-algebra. We can extend this argument to any collection of $\sigma$-algebras of $\Omega$, i.e., for any collection of $\sigma$-algebras, $\{B_j\}_{j\in\mathcal{J}}$, $\bigcap_{j\in\mathcal{J}}B_j$ is also a $\sigma$-algebra.
\end{proof}
\noindent Suppose $S$ is a collection of subsets of $\Omega$, which may or may not be a $\sigma$-algebra of $\Omega$. We want to embed this set $S$ into (possibly) larger collections, so that the new collection becomes a $\sigma$-algebra. Now, there exists a trivial $\sigma$-algebra, $2^\Omega$, which satisfies this requirement. However, there might be other smaller $\sigma$-algebras, $B\subseteq 2^\Omega$, such that $x\in S\implies x\in B$. In fact, it can be shown that there is a smallest of such $\sigma$-algebras.
\begin{theorem}
For every collection $S$ of subsets of $\Omega$, there exists a smallest $\sigma$-algebra containing all the elements of $S$, denoted by $\Sigma(S)$, i.e. for any other $\sigma$-algebra, $\Sigma'(S)$, that contains all the elements of $S$, we have $\Sigma(S)\subseteq\Sigma'(S)$.
\end{theorem}
\begin{proof}
Let us consider the collection of all $\sigma$-algebras, $\{B_j\}_{j\in\mathcal{J}}$, which contains all the elements of $S$. This collection is non-empty, because $2^\Omega\in \{B_j\}_{j\in\mathcal{J}}$. Now, if we take the intersection of all $B_j$, then, by theorem 4, the result, $\Sigma(S)=\bigcap_{j\in\mathcal{J}} B_j$, is a $\sigma$-algebra. Since we constructed this by taking the intersection of sets in $\{B_j\}_{j\in\mathcal{J}}$, it clearly contains all the elements of $S$. Now, let's consider another $\sigma$-algebra, $\Sigma'(S)$ which also contains all the elements of $S$. By design, $\Sigma'(S)\in \{B_j\}_{j\in\mathcal{J}}$. Therefore, $\Sigma(S)=\bigcap_{j\in\mathcal{J}} B_j\subseteq\Sigma'(S)$.
\end{proof}
\subsection{Borel $\sigma$-algebra and Borel Set}
Let's consider the collection of all open-intervals, $\mathcal{G}$, in a metric space, $(S, d)$. This collection, $\mathcal{G}$ is clearly not a $\sigma$-algebra, since the compliment of every open set in $\mathcal{G}$ is a close set by definition, and therefore doesn't belong in $\mathcal{G}$. However, by theorem 5, we can embed $\mathcal{G}$ to obtain a minimal $\sigma$-algebra that contains all the elements of $\mathcal{G}$. We can denote this as $\Sigma(\mathcal{G})$, known as the Borel $\sigma$-algebra.

In a similar fashion, we can obtain a Borel $\sigma$-algebra from the set of close sets as well, $\Sigma(\mathcal{F})$. We can show that these two constructions are identical.
\begin{theorem}
$\Sigma(\mathcal{G})=\Sigma(\mathcal{F})$.
\end{theorem}
\begin{proof}
$\forall x\in\mathcal{G}\implies x\in\Sigma(\mathcal{G})$. Since $\Sigma(\mathcal{G})$ is a $\sigma$-algebra, then, by definition, $x^c\in\Sigma(\mathcal{G})$. But the set of $x^c$ is $\mathcal{F}$. This means, that $\Sigma(\mathcal{G})$ contains all the elements of $\mathcal{F}$. So, from theorem 5, we have $\Sigma(\mathcal{G})\subseteq \Sigma(\mathcal{F})$.

Now, if we follow the same argument with $\forall y\in \mathcal{F}$, we see that $\Sigma(\mathcal{F})\subseteq \Sigma(\mathcal{G})$. Therefore, we can conclude that $\Sigma(\mathcal{G})=\Sigma(\mathcal{F})$.
\end{proof}
\noindent The sets that belong to $\Sigma(\mathcal{G})$ are called Borel sets. It is easy to see that the Borel $\sigma$-algebra include $\phi$, the entire metric space $S$, all the open-sets in $\mathcal{G}$, all the close sets in $\mathcal{F}$, and their countable union and intersections. We can also form higher form of Borel $\sigma$-algebras from $\mathcal{G}_\delta$, $\mathcal{F}_\sigma$ , $\mathcal{G}_{\delta\sigma}$, etc.
\end{document}